\section{Minimizing SSD-Level WAF}~\label{sec:ssdwa} % : SSD Hardware Interface-Agnostic Approaches

\module{SSDs also perform out-of-place writes.}
As discussed in \cref{sec:oop}, SSDs inherently use out-of-place writes and thus incur internal WA.
Along with the host write patterns, SSD WAF is shaped by the device’s internal architecture and algorithms~\cite{DBLP:conf/sigmod/LernerB21,lerner2024principles}.
Although SSD WAF is critical for performance and lifespan, it remains difficult to mitigate because SSDs operate as black boxes, exposing neither their internal behavior nor physical layout~\cite{ssdiq}.
Moreover, SSDs lack visibility into the workloads running above them.

\module{Different SSD interfaces for minimizing SSD WA.}
To bridge the knowledge gap between the host and SSD, recent work has explored more communicative SSD interfaces~\cite{DBLP:conf/sigmod/LernerB21}.
Zoned Namespace (ZNS) gives the host control over physical layout under certain write constraints, eliminating internal WA~\cite{DBLP:conf/hotstorage/Maheshwari21}.
Flexible Data Placement (FDP) allows the host to pass workload hints that the SSD can use to reduce internal WA~\cite{SSDhistory}.
While both interfaces are promising, their benefits depend on hardware support and proper integration with host-level management.

\module{Minimizing SSD WA regardless of SSD interfaces.}
We first show how our out-of-place write design naturally supports ZNS SSDs (\cref{sec:zns}).
When ZNS is unavailable and SSD-internal GC remains, 
we align DBMS and SSD GC units (\cref{sec:gcgranularity}) and introduce the \emph{NoWA write pattern}, which eliminates GC-induced WA even on commodity SSDs (\cref{sec:nowa}).
Finally, for FDP-enabled SSDs, we show how FDP placement hints can substitute for the NoWA pattern (\cref{sec:fdp}).

%%%%%%%%%%%%%%%%%%%%%%%%%%%%%%%%%%%%%%%%%%%%%%%%%%%%%%%%%%%%%%%%%%%%%%%%%%%%%%%%%%%%%%%%%%%%%%%%%%%%%%%%%%%%%%%%%%%%%%%%%
\subsection{Supporting Zoned Namespace SSD}~\label{sec:zns}

\module{SSD WAF = 1 with ZNS.}
A key benefit of ZNS SSDs is that they eliminate internal WA, making total WAF (DB WAF $\times$ SSD WAF) effectively equal to the DB WAF, as GC occurs only at the DB level.  
To guarantee SSD WAF of 1, ZNS enforces sequential writes within each zone and shifts GC responsibility to the host.
This fits naturally with our design that already uses a translation layer and out-of-place writes with its own GC~\cite{DBLP:conf/usenix/BjorlingAHRMGA21}.  
With full workload knowledge, GDT-based GC combined with ZNS can perform GC more effectively than in-place systems on commodity SSDs (\cref{fig:oop1}).

\module{ZNS provides more space and reduces DB WA.}
ZNS SSDs expose more usable capacity because they require no internal GC and therefore no OP space~\cite{DBLP:conf/usenix/BjorlingAHRMGA21}.  
Commodity SSDs typically reserve 7--28\% of capacity for OP~\cite{DBLP:conf/fast/ManeasMES22}, whereas ZNS returns this space to the host~\cite{DBLP:conf/fast/YadgarYS15}.  
This extra capacity reduces DB WA: from the DB GC’s perspective, it acts as additional OP space, allowing GC to be delayed until more pages become invalid.  
When GC finally runs, zones have lower valid ratios, directly reducing DB WA~\cite{arkiv/Dayan15}.

\module{Optimization: compatibility with ZNS.}
To align our layout with ZNS, we set the database zone size (GC unit) to the ZNS zone size.  
Each zone maintains minimal metadata, including a write pointer and its state (empty, open, full).  
Writes are issued using zone-append commands to zone IDs.  
To reuse a full zone after GC copyback, the host issues a zone reset~\cite{DBLP:conf/usenix/BjorlingAHRMGA21}.  
Overall, our design integrates cleanly with ZNS and remains straightforward to implement.
